% IEEE Conference Paper - KiwiDB Implementation
% Compile with: pdflatex kiwidb-implementation-paper.tex

\documentclass[conference]{IEEEtran}

\usepackage{amsmath,amssymb,amsfonts}
\usepackage{algorithmic}
\usepackage{graphicx}
\usepackage{textcomp}
\usepackage{xcolor}
\usepackage{listings}
\usepackage{booktabs}
\usepackage{multirow}
\usepackage{url}
\usepackage{hyperref}
\usepackage{float}
\usepackage{enumitem}

\lstset{
  basicstyle=\footnotesize\ttfamily,
  breaklines=true,
  frame=single,
  numbers=left,
  numberstyle=\tiny,
  tabsize=2,
  captionpos=b
}

\begin{document}

\title{KiwiDB: An Async-Native Log-Structured Merge Tree Key-Value Store with Parallel Flush and Subprocess Compaction}

\author{
\IEEEauthorblockN{Chaitanya Paraskar}
\IEEEauthorblockA{Department of Computer Engineering \\
University of Pune \\
Pune, India \\
chaitanya.paraskar@example.com}
% Add additional authors as needed
}

\maketitle

\begin{abstract}
Log-Structured Merge Trees (LSM trees) are the dominant write-optimized storage engine architecture, underpinning systems such as LevelDB, RocksDB, and Apache Cassandra. However, existing implementations are written in C/C++ with synchronous APIs, making them incompatible with modern asynchronous application frameworks. This paper presents KiwiDB, a production-grade LSM tree key-value store implemented from scratch in Python 3.12+. KiwiDB introduces three key design innovations: (1) an async-native API where the event loop never blocks on disk I/O; (2) a parallel flush pipeline with event-chain ordered commits that achieves concurrent SSTable writes while preserving strict FIFO ordering; and (3) subprocess-based compaction that bypasses CPython's Global Interpreter Lock (GIL) for true write-path isolation. The engine employs a hybrid L0-tiering/L1-leveling compaction strategy that reduces write amplification by an order of magnitude compared to pure leveling, a concurrent skip list with lock-free reads and per-node write locking, data-derived bloom filter sizing, and a three-tier block cache. KiwiDB is accompanied by a comprehensive observability layer including a FastAPI REST interface and a React-based web dashboard for real-time engine inspection. The system comprises approximately 6,600 lines of type-checked Python with 211 tests across 25 modules, demonstrating that a well-engineered LSM engine can be built from first principles with formal type safety and modern concurrency primitives.
\end{abstract}

\begin{IEEEkeywords}
LSM tree, key-value store, async I/O, compaction, skip list, bloom filter, write-ahead log, Python, storage engine
\end{IEEEkeywords}

% ============================================================
\section{Introduction}
% ============================================================

The Log-Structured Merge Tree (LSM tree), introduced by O'Neil et al.~\cite{oneil1996}, has become the de facto storage engine architecture for write-heavy workloads. By buffering writes in memory and flushing them as immutable sorted runs to disk, LSM trees convert random writes into sequential I/O --- achieving write throughput that B-tree based systems cannot match. This design powers production systems including Google's LevelDB~\cite{leveldb}, Facebook's RocksDB~\cite{rocksdb}, Apache Cassandra~\cite{cassandra}, and ScyllaDB.

Despite the maturity of these systems, they share two limitations that motivated this work:

\textbf{Synchronous APIs.} LevelDB and RocksDB expose blocking C++ interfaces. Integrating them into modern async frameworks (FastAPI, aiohttp, Node.js) requires external thread-pool wrappers that add latency and complexity. No production LSM implementation offers a natively asynchronous API.

\textbf{Opaque internals.} Production engines prioritize throughput over observability. Inspecting memtable state, bloom filter hit rates, compaction progress, or SSTable layouts requires external tooling (e.g., \texttt{sst\_dump}, Grafana dashboards). The internal design decisions are documented sparsely, making these systems difficult to learn from.

KiwiDB addresses both limitations by implementing a complete LSM tree from scratch in Python 3.12+ with the following contributions:

\begin{enumerate}[leftmargin=*]
    \item \textbf{Async-native API}: All public operations (\texttt{put}, \texttt{get}, \texttt{delete}, \texttt{flush}, \texttt{close}) are async coroutines. WAL fsync, SSTable flush, and compaction execute off the event loop via \texttt{asyncio.to\_thread()} and \texttt{ProcessPoolExecutor}.

    \item \textbf{Parallel flush with event-chain ordering}: Multiple memtables are flushed to SSTables concurrently, but commits are serialized in oldest-first order via an \texttt{asyncio.Event} chain --- doubling flush throughput without sacrificing correctness.

    \item \textbf{Subprocess compaction with GIL bypass}: Compaction merges execute in separate processes via \texttt{ProcessPoolExecutor}, achieving true CPU parallelism that is impossible with Python threads due to the GIL.

    \item \textbf{Hybrid compaction strategy}: L0 uses tiering (multiple overlapping files) while L1--L3 use leveling (one file per level), reducing write amplification by $10\times$ compared to LevelDB's pure leveling.

    \item \textbf{Built-in observability}: A FastAPI REST interface and React web dashboard provide real-time inspection of memtables, SSTables, compaction jobs, and engine logs via WebSocket streaming.
\end{enumerate}

The remainder of this paper is organized as follows. Section~\ref{sec:related} surveys related work. Section~\ref{sec:architecture} presents the system architecture and layered dependency model. Section~\ref{sec:design} details the design of each subsystem. Section~\ref{sec:implementation} covers implementation-specific decisions, concurrency techniques, and Python-specific challenges. Section~\ref{sec:conclusion} concludes with a discussion of limitations and future work.

% ============================================================
\section{Related Work}
\label{sec:related}
% ============================================================

\subsection{LSM Tree Foundations}

The LSM tree was proposed by O'Neil et al.~\cite{oneil1996} as a write-optimized alternative to B-trees. The key insight is that buffering writes in memory and flushing them as sorted runs to disk converts random writes into sequential I/O. Subsequent work by Luo and Carey~\cite{luo2020} formalized the design space across three dimensions: tiering vs.\ leveling, size ratio, and merge policy.

\subsection{Production LSM Implementations}

\textbf{LevelDB}~\cite{leveldb} by Google is the canonical LSM implementation. It uses pure leveling compaction, a single-threaded compaction model, and a skip list memtable protected by a global mutex. Its API is synchronous C++.

\textbf{RocksDB}~\cite{rocksdb} by Facebook extends LevelDB with concurrent compaction threads, column families, configurable compaction strategies, and ACID transactions. It remains synchronous with optional async interfaces introduced in version 8.x.

\textbf{WiredTiger}~\cite{wiredtiger}, the MongoDB storage engine, combines B-tree and LSM approaches with multiversion concurrency control (MVCC). It is implemented in C with a synchronous API.

\subsection{LSM Optimizations}

Bloom filters~\cite{bloom1970} are standard for reducing read amplification in LSM trees by avoiding unnecessary disk reads. Monkey~\cite{dayan2017} optimizes bloom filter memory allocation across levels by assigning more bits to deeper levels. KiwiDB takes a different approach: it sizes each filter using the actual record count at write time, producing optimally sized filters without cross-level optimization.

Write amplification in leveled compaction has been extensively studied. Dostoevsky~\cite{dayan2018} introduces lazy leveling as a hybrid between tiering and leveling. KiwiDB's hybrid strategy --- L0 tiering with L1+ leveling --- achieves a similar write-amplification reduction through a simpler design.

\subsection{Concurrent Data Structures}

The concurrent skip list in KiwiDB draws on the lock-free skip list by Herlihy et al.~\cite{herlihy2006}. However, KiwiDB adapts the design for CPython's GIL: writes use per-node fine-grained locking (rather than compare-and-swap), while reads are lock-free by relying on GIL-atomic boolean flag reads.

% ============================================================
\section{System Architecture}
\label{sec:architecture}
% ============================================================

\subsection{Overview}

KiwiDB is structured as a layered system with strict dependency ordering. The engine coordinates six subsystems: Write-Ahead Log (WAL), MemTable, SSTable, Flush Pipeline, Compaction, and Observability. No layer imports from a layer above it, preventing circular dependencies and enabling independent testing at each tier.

\subsection{Layered Dependency Model}

The codebase is organized into seven tiers, from leaf modules with no internal dependencies up to the top-level engine coordinator:

\begin{table}[H]
\centering
\caption{KiwiDB Layered Dependency Model}
\label{tab:layers}
\begin{tabular}{@{}clc@{}}
\toprule
\textbf{Tier} & \textbf{Modules} & \textbf{Depends On} \\
\midrule
0 & \texttt{types}, \texttt{crc}, \texttt{uuid7}, \texttt{errors} & stdlib only \\
1 & \texttt{bloom}, \texttt{cache}, \texttt{index}, \texttt{rwlock} & Tier 0 \\
2 & \texttt{memtable}, \texttt{wal} & Tiers 0--1 \\
3 & \texttt{sstable}, \texttt{config}, \texttt{seq\_gen} & Tiers 0--2 \\
4 & \texttt{memtable\_mgr}, \texttt{sstable\_mgr}, \texttt{compact/} & Tiers 0--3 \\
5 & \texttt{flush\_pipeline}, \texttt{compaction\_mgr} & Tiers 0--4 \\
6 & \texttt{lsm\_engine} (public API) & All tiers \\
\bottomrule
\end{tabular}
\end{table}

This strict ordering ensures that each tier can be unit-tested in isolation. The engine at Tier 6 is a thin coordinator that delegates all state management to the subsystem managers below it.

\subsection{Write Path}

Every write operation follows an atomic four-step sequence under a single lock, enforcing the invariant \textit{durability before visibility}:

\begin{enumerate}[leftmargin=*]
    \item \textbf{Sequence generation}: \texttt{SeqGenerator.next()} atomically increments a monotonic counter under \texttt{threading.Lock}. The sequence number serves as the MVCC version --- higher values indicate newer data.

    \item \textbf{WAL durability}: The entry is serialized (msgpack) with a CRC32 checksum and appended to the WAL file. The file descriptor is fsynced before returning, ensuring the write survives process or OS crashes.

    \item \textbf{MemTable insert}: The key-value pair is inserted into the active skip list via \texttt{MemTableManager.put()}.

    \item \textbf{Conditional freeze}: If the active memtable exceeds its threshold (entry count in dev mode, byte size in production), it is frozen into an \texttt{ImmutableMemTable}, pushed to the immutable queue, and the flush pipeline is signaled.
\end{enumerate}

All four steps execute under \texttt{\_mem.write\_lock}. Either all succeed or the caller observes the exception --- partial writes are impossible.

\subsection{Read Path}

Reads scan data sources from newest to oldest, returning the first match with the highest sequence number:

\begin{enumerate}[leftmargin=*]
    \item \textbf{Active memtable}: $O(\log n)$ lock-free skip list traversal.
    \item \textbf{Immutable queue}: Newest-first scan of 0--4 frozen snapshots. Each lookup is $O(1)$ via an internal dictionary index.
    \item \textbf{L0 SSTables}: All files are checked (key ranges overlap). For each file: bloom filter $\rightarrow$ sparse index bisect $\rightarrow$ block scan. The result with the highest sequence number across all L0 files is kept.
    \item \textbf{L1--L3 SSTables}: One file per level, checked sequentially. The same bloom $\rightarrow$ index $\rightarrow$ block pipeline applies.
\end{enumerate}

Tombstones (deletion markers) are detected at the final step: if the matched value equals the sentinel \texttt{b"\textbackslash x00\_\_tomb\_\_\textbackslash x00"}, the engine returns \texttt{None}.

\subsection{Recovery Model}

On startup, KiwiDB reconstructs in-memory state from durable storage:

\begin{enumerate}[leftmargin=*]
    \item Load configuration and open WAL file.
    \item Load all SSTables from disk via the persistent manifest.
    \item Compute \texttt{max\_seq\_seen} across all loaded SSTables.
    \item Replay WAL entries where \texttt{seq > max\_seq\_seen} into a fresh memtable.
    \item Restore the sequence generator to the highest observed sequence.
    \item Start the flush pipeline and compaction manager.
\end{enumerate}

Entries already flushed to SSTables are skipped during replay, ensuring no duplicates. The completeness of an SSTable is signaled by the presence of its \texttt{meta.json} file --- incomplete SSTables (from crash during write) are ignored.

% ============================================================
\section{Subsystem Design}
\label{sec:design}
% ============================================================

\subsection{Write-Ahead Log (WAL)}

The WAL provides crash recovery by persisting every write before it becomes visible in the memtable. Each entry is serialized as a msgpack-framed record with CRC32 integrity verification:

\begin{center}
\texttt{[4B length][msgpack(seq, ts, op, key, val)][4B CRC32]}
\end{center}

Two operation types are supported: \texttt{PUT} (op=1) and \texttt{DELETE} (op=2). Deletes are written as tombstone entries that propagate through flush and compaction.

\textbf{Durability guarantee}: Every append performs \texttt{os.fsync()} on the file descriptor before returning. This ensures that even if the process crashes immediately after \texttt{put()} returns, the entry is recoverable from disk.

\textbf{Lock ordering}: The WAL lock (\texttt{\_wal\_lock}) is always acquired \emph{after} the memtable write lock (\texttt{\_mem.write\_lock}). Reversing this order would cause deadlock, since the write path holds the memtable lock while calling \texttt{sync\_append()}, which internally acquires the WAL lock.

\textbf{Truncation}: After a successful flush commit, the WAL is truncated to remove entries that are now durable in SSTables. This prevents unbounded WAL growth.

\subsection{MemTable and Concurrent Skip List}

\subsubsection{Skip List Data Structure}

The memtable is backed by a concurrent skip list --- a probabilistic sorted data structure that provides $O(\log n)$ expected time for insertions and lookups. KiwiDB's skip list uses a maximum of 16 levels with a promotion probability of 0.5 (geometric distribution), supporting approximately 65,536 entries before the highest level is reached.

Each node stores:
\begin{itemize}[leftmargin=*]
    \item \texttt{key: bytes} --- the lookup key
    \item \texttt{seq: int} --- MVCC sequence number
    \item \texttt{timestamp\_ms: int} --- insertion timestamp
    \item \texttt{value: bytes} --- the stored value (or tombstone sentinel)
    \item \texttt{forward: list[Node]} --- forward pointers (one per level)
    \item \texttt{lock: threading.Lock} --- per-node fine-grained lock
    \item \texttt{marked: bool} --- logical deletion flag
    \item \texttt{fully\_linked: bool} --- visibility flag
\end{itemize}

\subsubsection{Concurrency Model}

The skip list implements a fine-grained locking protocol adapted for CPython's GIL:

\textbf{Writers} lock only the predecessor nodes at the insertion boundary. Two concurrent inserts with disjoint keys do not contend. The insertion protocol is:
\begin{enumerate}[leftmargin=*]
    \item Top-down traversal to find predecessor and successor at each level.
    \item Lock predecessor nodes at the insertion boundary.
    \item Validate that predecessors are still correct (not marked or unlinked).
    \item Create the new node, link it at level 0 (becomes visible to readers).
    \item Link at higher levels, then set \texttt{fully\_linked = True}.
    \item Release all locks.
\end{enumerate}

\textbf{Readers} perform a lock-free top-down traversal, checking the \texttt{fully\_linked} and \texttt{marked} boolean flags. Under CPython, single-attribute reads on Python objects are GIL-atomic --- no torn reads are possible. This makes reader traversal effectively lock-free without requiring explicit compare-and-swap (CAS) instructions.

\subsubsection{MemTable Manager}

The \texttt{MemTableManager} coordinates the active memtable, immutable queue, and freeze signaling:

\begin{itemize}[leftmargin=*]
    \item \textbf{Active memtable}: The currently writable skip list instance.
    \item \textbf{Immutable queue}: A bounded deque of up to 4 frozen snapshots awaiting flush. When full, new writes block until the flush pipeline drains space (backpressure).
    \item \textbf{Freeze}: When the active memtable exceeds its threshold, it is snapshot into an \texttt{ImmutableMemTable} (sorted list + dictionary index for $O(1)$ lookups), pushed to the queue, and a fresh active memtable is created.
    \item \textbf{Signal bridging}: The freeze event is a \texttt{threading.Event} bridged to async via \texttt{loop.call\_soon\_threadsafe()}, allowing the flush pipeline (async) to wake up immediately when a sync write triggers a freeze.
\end{itemize}

\subsection{SSTable Format}

Each SSTable is stored as a four-file directory:

\begin{table}[H]
\centering
\caption{SSTable Directory Layout}
\label{tab:sstable}
\begin{tabular}{@{}ll@{}}
\toprule
\textbf{File} & \textbf{Purpose} \\
\midrule
\texttt{data.bin} & Sorted KV records in $\sim$4\,KB blocks \\
\texttt{index.bin} & Sparse index (first key per block $\rightarrow$ byte offset) \\
\texttt{filter.bin} & Bloom filter bit array + header \\
\texttt{meta.json} & Metadata (written last as completeness signal) \\
\bottomrule
\end{tabular}
\end{table}

\subsubsection{Record Encoding}

Each record in \texttt{data.bin} uses a fixed-header binary format:

\begin{center}
\texttt{[2B key\_len][2B val\_len][8B seq][8B ts\_ms][key][value][4B CRC32]}
\end{center}

Key and value lengths are 2-byte big-endian unsigned integers. Sequence and timestamp are 8-byte big-endian signed integers. The CRC32 covers the entire record (header + payload) for integrity verification.

Records are grouped into blocks of approximately 4,096 bytes. When the write buffer exceeds the block size threshold, it is flushed to disk. The sparse index records the first key and byte offset of each flushed block.

\subsubsection{Writer State Machine}

The SSTable writer follows a strict write-once lifecycle:

\begin{center}
\texttt{OPEN $\rightarrow$ put() $\times$ N $\rightarrow$ finish() $\rightarrow$ DONE}
\end{center}

Keys must be inserted in ascending order (enforced by the writer). The \texttt{finish()} method is async for L0 flushes (writing bloom filter and index concurrently via \texttt{asyncio.gather}), while \texttt{finish\_sync()} is synchronous for compaction subprocesses. The \texttt{meta.json} file is written last --- its presence signals that the SSTable is complete and safe to read.

\subsubsection{Reader and Memory-Mapped I/O}

The SSTable reader uses memory-mapped I/O (\texttt{mmap}) for zero-copy access to \texttt{data.bin}. The bloom filter and sparse index are loaded lazily on the first \texttt{get()} call and cached in the shared block cache for cross-restart reuse.

Point lookup follows this pipeline:
\begin{enumerate}[leftmargin=*]
    \item \textbf{Bloom filter check}: If the key is definitely absent, skip this SSTable entirely.
    \item \textbf{Sparse index bisect}: \texttt{bisect\_right} finds the block whose first key $\leq$ the search key.
    \item \textbf{Block scan}: The identified block is read from cache (or mmap on miss) and scanned sequentially. Since records are sorted, the scan terminates early when the current key exceeds the search key.
\end{enumerate}

\subsection{Bloom Filter}

KiwiDB's bloom filter uses MurmurHash3 (via the \texttt{mmh3} library) with multiple hash functions derived from different seeds. The filter is sized using the optimal formulas:

\begin{equation}
m = \left\lceil \frac{-n \cdot \ln p}{(\ln 2)^2} \right\rceil
\label{eq:bloom_bits}
\end{equation}

\begin{equation}
k = \left\lceil \frac{m}{n} \cdot \ln 2 \right\rceil
\label{eq:bloom_hashes}
\end{equation}

where $n$ is the number of items, $p$ is the target false positive rate, $m$ is the bit array size, and $k$ is the number of hash functions.

\textbf{Data-derived sizing}: Unlike LevelDB and RocksDB which use a fixed bits-per-key ratio, KiwiDB derives $n$ from the actual record count at write time --- \texttt{len(snapshot)} for flush, sum of input record counts for compaction. This produces an optimally sized filter for every SSTable without manual tuning.

\textbf{Environment-aware FPR}: The target false positive rate is configurable per environment: 5\% in development (faster builds, larger filters) and 1\% in production (fewer false disk reads).

The filter is serialized as a binary header (number of hashes, bit count, seed) followed by the byte array and a CRC32 checksum.

\subsection{Sparse Index}

The sparse index maps the first key of each data block to its byte offset in \texttt{data.bin}. It supports two key operations:

\begin{itemize}[leftmargin=*]
    \item \texttt{floor\_offset(key)}: Returns the offset of the block whose first key is $\leq$ the search key, using \texttt{bisect\_right}.
    \item \texttt{ceil\_offset(key)}: Returns the offset of the block whose first key is $\geq$ the search key, using \texttt{bisect\_left}.
\end{itemize}

The index is serialized as variable-length entries (\texttt{[4B key\_len][8B offset][key\_bytes]}) with a trailing CRC32. At query time, the index is loaded into memory and searched via Python's \texttt{bisect} module, providing $O(\log B)$ lookup time where $B$ is the number of blocks.

\subsection{Block Cache}

KiwiDB implements a three-tier LRU cache that serves data blocks, sparse indexes, and bloom filters with independent eviction policies:

\begin{table}[H]
\centering
\caption{Block Cache Tiers}
\label{tab:cache}
\begin{tabular}{@{}lccc@{}}
\toprule
\textbf{Tier} & \textbf{Offset Key} & \textbf{Default Cap.} & \textbf{Eviction Priority} \\
\midrule
Data blocks & $\geq 0$ & 256 entries & Lowest (evicted first) \\
Sparse indexes & $-2$ & 64 entries & Medium \\
Bloom filters & $-1$ & 64 entries & Highest (evicted last) \\
\bottomrule
\end{tabular}
\end{table}

The tiered design reflects access patterns: bloom filters are checked on every lookup and should remain resident, indexes are checked on positive bloom results, and data blocks are accessed only when the index indicates a candidate block. Each tier uses an independent LRU backed by \texttt{cachetools.LRUCache} with \texttt{threading.RLock} for thread safety.

\subsection{Flush Pipeline}

The flush pipeline is an async daemon that drains the immutable queue to disk. Its key innovation is \textit{parallel writes with ordered commits}.

\subsubsection{Parallel Writes}

Up to \texttt{flush\_max\_workers} (default: 2) SSTable writes execute concurrently, bounded by an \texttt{asyncio.Semaphore}. Each write targets a separate directory, so there are no I/O conflicts.

\subsubsection{Event-Chain Ordered Commits}

Writes are parallel, but commits must occur in oldest-first order to maintain the invariant that L0 files are ordered by age. This is achieved through an \texttt{asyncio.Event} chain:

Each flush slot carries two events:
\begin{itemize}[leftmargin=*]
    \item \texttt{prev\_committed}: Set by the predecessor slot upon commit.
    \item \texttt{my\_committed}: Set by this slot upon its own commit.
\end{itemize}

After completing its disk write, each worker blocks on \texttt{prev\_committed.wait()} before committing. The wait does not delay the write --- only the commit. This achieves maximum write parallelism while serializing commits.

\subsubsection{Correctness Invariant}

The core invariant is: \textit{at every point in time, every durable key must be findable by \texttt{get()} --- either in the immutable queue or in the SSTable registry, never in neither.}

This is maintained because:
\begin{enumerate}[leftmargin=*]
    \item A snapshot stays in the immutable queue until its commit completes.
    \item The commit atomically registers the new SSTable reader \emph{before} popping the snapshot from the queue.
    \item WAL truncation happens \emph{after} both registration and pop.
\end{enumerate}

\subsubsection{Crash Safety}

The pipeline is crash-safe at every point:
\begin{itemize}[leftmargin=*]
    \item \textbf{Crash during write}: Snapshot remains in immutable queue. Incomplete SSTable (no \texttt{meta.json}) is ignored on recovery. WAL replays all entries.
    \item \textbf{Crash after register, before pop}: SSTable is readable; snapshot is redundant but safe. WAL replay deduplicates by sequence number.
    \item \textbf{Crash after pop, before WAL truncate}: SSTable has all data. WAL entries are replayed and deduplicated.
    \item \textbf{Crash after WAL truncate}: Clean state --- SSTable has data, WAL is trimmed.
\end{itemize}

\subsubsection{Failure Propagation}

If any slot fails during the write phase:
\begin{enumerate}[leftmargin=*]
    \item A \texttt{batch\_abort} event is set, causing all downstream slots to skip their commits.
    \item \texttt{my\_committed} is set in the \texttt{finally} block, preventing chain deadlock.
    \item The failed snapshot remains in the queue for retry on the next dispatch cycle.
    \item The WAL is not truncated, preserving data recoverability.
\end{enumerate}

\subsection{Compaction}

\subsubsection{Level Model}

KiwiDB uses a hybrid compaction strategy combining L0 tiering with L1+ leveling:

\begin{itemize}[leftmargin=*]
    \item \textbf{L0 (Tiering)}: Up to 10 files with overlapping key ranges. Flushed directly from memtable with no merge --- a pure sequential append.
    \item \textbf{L1--L3 (Leveling)}: Exactly one file per level, produced by merging the level above with the existing file at the target level.
\end{itemize}

\subsubsection{Trigger Mechanism}

Compaction is flush-triggered, not daemon-driven. After every successful flush commit, the flush pipeline calls \texttt{CompactionManager.check\_and\_compact()}. This avoids idle polling and ensures compaction runs promptly.

\begin{table}[H]
\centering
\caption{Compaction Trigger Thresholds}
\label{tab:compact}
\begin{tabular}{@{}lll@{}}
\toprule
\textbf{Transition} & \textbf{Trigger} & \textbf{Config Key} \\
\midrule
L0 $\rightarrow$ L1 & File count $\geq$ threshold & \texttt{l0\_compaction\_threshold} \\
L1 $\rightarrow$ L2 & Size exceeds $10^2 \times$ base & Cascading formula \\
L2 $\rightarrow$ L3 & Size exceeds $10^3 \times$ base & Cascading formula \\
\bottomrule
\end{tabular}
\end{table}

\subsubsection{K-Way Merge Algorithm}

The merge uses a \texttt{KWayMergeIterator} --- a min-heap over $N$ sorted input iterators. The comparison key is \texttt{(key, -seq)}: keys are sorted ascending, and for equal keys, the highest sequence number (newest write) is sorted first. This enables single-pass deduplication: for each unique key, only the first (highest-seq) record is emitted.

\textbf{Tombstone garbage collection}: Tombstones with \texttt{seq < seq\_cutoff} are dropped during merge, as the deletion has propagated to all levels. Tombstones above the cutoff are preserved to shadow values in deeper levels that have not yet been compacted.

\subsubsection{Subprocess Execution}

Each compaction job is represented by a \texttt{CompactionTask} --- a frozen dataclass with only primitive fields (strings, integers, lists of strings) that can safely cross the subprocess boundary. The merge executes in a \texttt{ProcessPoolExecutor} subprocess with its own GIL, memory space, and file handles. The main process's event loop remains free throughout.

\subsubsection{Atomic Commit}

After the subprocess produces a new SSTable, the commit is atomic:
\begin{enumerate}[leftmargin=*]
    \item Register new reader (destination level becomes readable).
    \item Write manifest to disk (durable).
    \item Mark old files for deletion (deferred by reference count).
    \item Update in-memory state.
    \item Evict stale cache blocks.
\end{enumerate}

Write locks on source and destination levels prevent reads from seeing inconsistent state during the brief commit phase.

\subsection{SSTable Manager and Manifest}

The \texttt{SSTableManager} coordinates all on-disk state via a persistent JSON manifest:

\begin{lstlisting}[caption={Manifest structure},label={lst:manifest}]
{
  "l0_order": ["file_a", "file_b"],
  "levels": {"1": "file_c", "2": "file_d"}
}
\end{lstlisting}

The \texttt{l0\_order} list stores file IDs in newest-first order, ensuring reads check the most recent data first. Per-level \texttt{AsyncRWLock} instances serialize compaction commits while allowing concurrent readers. A reference-counted registry manages SSTable reader lifecycle, preventing reader closure while in-flight reads are active.

% ============================================================
\section{Implementation Details}
\label{sec:implementation}
% ============================================================

\subsection{Async/Sync Bridge}

KiwiDB exposes an async-first API, but several low-level operations are inherently synchronous (memtable inserts, WAL fsync, skip list traversal). The bridge between these worlds uses three patterns:

\textbf{Pattern 1 --- Direct sync under lock}: Operations that complete in microseconds (sequence generation, memtable insert) execute synchronously within an async method. The caller awaits the method, but the actual work is a fast synchronous critical section.

\textbf{Pattern 2 --- \texttt{asyncio.to\_thread()}}: Blocking I/O operations (WAL fsync, SSTable reads via mmap) are offloaded to the default thread pool, freeing the event loop.

\textbf{Pattern 3 --- \texttt{ProcessPoolExecutor}}: CPU-bound operations (compaction merge) are dispatched to subprocesses. The dispatch itself is wrapped in \texttt{asyncio.to\_thread(pool.submit(...).result)} to avoid blocking the event loop during IPC.

\subsection{GIL Bypass for Compaction}

CPython's Global Interpreter Lock (GIL) prevents true parallelism between threads in a single process. This is acceptable for I/O-bound operations (where threads release the GIL during syscalls) but problematic for CPU-bound compaction merges that must iterate millions of records.

KiwiDB solves this by running compaction in a separate process via \texttt{ProcessPoolExecutor}. The subprocess has its own Python interpreter and GIL, achieving true parallelism. The trade-off is IPC overhead ($\sim$5--10\,ms for process spawn and result serialization), which is negligible compared to the multi-second merge duration.

The \texttt{CompactionTask} dataclass is designed for subprocess compatibility: all fields are primitive types (strings, integers, lists of strings) with no references to open file handles, locks, or asyncio objects.

\subsection{Lock Ordering and Deadlock Prevention}

KiwiDB enforces a strict lock ordering discipline to prevent deadlocks:

\begin{center}
\texttt{\_mem.write\_lock} $\rightarrow$ \texttt{\_wal\_lock} $\rightarrow$ \texttt{level[i].rwlock}
\end{center}

The write path always acquires \texttt{\_mem.write\_lock} first, then calls \texttt{sync\_append()} which acquires \texttt{\_wal\_lock}. Level-specific read/write locks are acquired independently during reads and compaction commits. No code path acquires these locks in reverse order.

\subsection{Writer-Preferred AsyncRWLock}

KiwiDB implements a custom \texttt{AsyncRWLock} with writer preference: once a writer is waiting, new readers block until the writer completes. This prevents writer starvation under continuous read load --- critical for compaction commits, which must acquire write locks on affected levels.

The lock maintains three counters: active readers, active writers (0 or 1), and waiting writers. A reader is admitted only if there are no active or waiting writers. This design ensures compaction commits are not indefinitely deferred by a stream of reads.

\subsection{Backpressure Mechanism}

When the immutable queue reaches its maximum length (default: 4), new writes block until the flush pipeline drains at least one snapshot. This backpressure mechanism prevents unbounded memory growth when writes outpace flush throughput.

The blocking is implemented via a \texttt{threading.Event} with a 60-second timeout. If the timeout expires, a \texttt{FreezeBackpressureTimeout} exception is raised, indicating that the flush pipeline may be stuck.

\subsection{Signal Bridging: Sync to Async}

A key implementation challenge is signaling between synchronous write operations and the asynchronous flush pipeline. When a sync write triggers a memtable freeze, the flush pipeline (running as an async daemon task) must wake up immediately.

KiwiDB bridges this via \texttt{loop.call\_soon\_threadsafe()}: the synchronous freeze code obtains a reference to the running event loop and schedules an \texttt{asyncio.Event.set()} call that wakes the flush pipeline on the next event loop iteration. This avoids polling and provides sub-millisecond wake-up latency.

\subsection{Thread and Async Model Summary}

Table~\ref{tab:threading} summarizes the concurrency model chosen for each component and the rationale behind each choice.

\begin{table}[H]
\centering
\caption{Thread and Async Model by Component}
\label{tab:threading}
\begin{tabular}{@{}lll@{}}
\toprule
\textbf{Component} & \textbf{Model} & \textbf{Rationale} \\
\midrule
MemTable ops & Sync (threading) & In-memory, $\mu$s latency \\
WAL append & Sync + \texttt{to\_thread} & Must fsync before return \\
Flush pipeline & Async daemon & Non-blocking parallel I/O \\
Compaction & Subprocess & CPU-bound, GIL bypass \\
SSTable reads & Sync (mmap) & Zero-copy after page warm \\
Config access & Sync (RLock) & Thread-safe attributes \\
\bottomrule
\end{tabular}
\end{table}

\subsection{Runtime-Mutable Configuration}

All engine configuration parameters (memtable thresholds, bloom filter FPR, compaction thresholds, worker counts) are stored in a thread-safe \texttt{Config} object backed by a JSON file on disk.

Configuration values are read at the point of use, not cached. Changing a parameter via \texttt{update\_config()} takes effect on the next operation that reads it --- no engine restart is required. Writes to the config file use atomic replacement (\texttt{tempfile} + \texttt{os.replace()}) to prevent corruption from concurrent writes or crashes.

\begin{table}[H]
\centering
\caption{Key Configuration Parameters}
\label{tab:config}
\begin{tabular}{@{}lccp{3.5cm}@{}}
\toprule
\textbf{Parameter} & \textbf{Dev} & \textbf{Prod} & \textbf{Purpose} \\
\midrule
\texttt{max\_memtable\_entries} & 10 & --- & Freeze trigger (count) \\
\texttt{max\_memtable\_size\_mb} & --- & 64 & Freeze trigger (bytes) \\
\texttt{l0\_compact\_threshold} & 10 & 10 & L0 file count limit \\
\texttt{bloom\_fpr} & 0.05 & 0.01 & Target false positive rate \\
\texttt{flush\_max\_workers} & 2 & 2 & Parallel flush workers \\
\texttt{block\_size} & 4096 & 4096 & SSTable block size (B) \\
\texttt{max\_levels} & 3 & 3 & Compaction depth \\
\bottomrule
\end{tabular}
\end{table}

\subsection{Type Safety and Formal Contracts}

KiwiDB enforces type safety through two independent static type checkers running in strict mode:

\begin{itemize}[leftmargin=*]
    \item \textbf{mypy} (strict mode): Disallows implicit \texttt{Any}, requires full annotations on all public APIs.
    \item \textbf{basedpyright} (strict mode): Provides additional type narrowing and flow analysis beyond mypy's capabilities.
\end{itemize}

Formal contracts are defined through:
\begin{itemize}[leftmargin=*]
    \item \textbf{Abstract Base Classes (ABCs)}: \texttt{StorageEngine} (defines the public API contract), \texttt{Serializable} (defines serialization/deserialization), \texttt{MemTable} (defines memtable operations).
    \item \textbf{Protocols}: Structural subtyping for interfaces that don't require inheritance.
    \item \textbf{.pyi stub files}: Type stubs for public API, enabling IDE autocompletion and third-party type checking.
\end{itemize}

\subsection{Observability Layer}

KiwiDB includes a comprehensive observability layer built on three components:

\textbf{Structured logging}: All engine events are logged via \texttt{structlog} with structured key-value fields (sequence numbers, file IDs, level numbers, byte sizes). Logs are written to both a rotating file and a TCP broadcast server for live consumption.

\textbf{REST API}: A FastAPI application exposes 20+ endpoints organized into seven categories: key-value operations (\texttt{GET/POST/DELETE /kv/\{key\}}), memtable inspection (\texttt{GET /mem}), SSTable browsing (\texttt{GET /disk}), engine statistics (\texttt{GET /stats}), configuration management (\texttt{PATCH /config}), engine control (\texttt{POST /engine/reset}), and a terminal interface (\texttt{POST /terminal/command}).

\textbf{Web dashboard}: A React single-page application (built with Vite and TypeScript) provides real-time visualization of memtable fill level and freeze events, SSTable layout by level, compaction job progress, live engine log stream via WebSocket, and an interactive browser-based REPL.

% ============================================================
\section{Comparison with Existing Systems}
\label{sec:comparison}
% ============================================================

Table~\ref{tab:comparison} summarizes the key architectural differences between KiwiDB, LevelDB, and RocksDB.

\begin{table*}[t]
\centering
\caption{Architectural Comparison: KiwiDB vs.\ LevelDB vs.\ RocksDB}
\label{tab:comparison}
\begin{tabular}{@{}llll@{}}
\toprule
\textbf{Dimension} & \textbf{KiwiDB} & \textbf{LevelDB} & \textbf{RocksDB} \\
\midrule
Language & Python 3.12+ & C++ & C++ \\
API model & Async-first (native \texttt{asyncio}) & Synchronous & Synchronous (partial async in 8.x) \\
L0 strategy & Tiering (append, no merge) & Leveling (merge) & Configurable \\
Flush pipeline & Parallel writes + ordered commits & Sequential (single thread) & Concurrent \\
Compaction isolation & Subprocess (GIL bypass) & Background thread & Thread pool \\
Skip list reads & Lock-free (GIL-atomic flags) & Mutex-locked & Lock-free (CAS) \\
Bloom filter sizing & Data-derived (optimal per SSTable) & Fixed bits-per-key & Fixed bits-per-key \\
Configuration & Full runtime mutability & Fixed at open & Partial mutability \\
Type safety & Dual strict checkers (mypy + pyright) & Compiler only & Compiler only \\
Built-in dashboard & React + FastAPI + WebSocket & None & None \\
Documentation & 15+ design docs + autodoc API ref & 2 docs & Wiki \\
\bottomrule
\end{tabular}
\end{table*}

\textbf{Write amplification analysis}: KiwiDB's L0 tiering eliminates the merge cost at flush time. A flush in KiwiDB costs $O(L/B)$ (sequential append of $L$-byte records in $B$-byte blocks), while LevelDB's pure leveling costs $O(T \times L/B)$ where $T$ is the size ratio --- a $T\times$ reduction (approximately $10\times$ with the default threshold of 10).

\textbf{Read amplification trade-off}: KiwiDB checks up to $T$ bloom filters at L0 (one per file), plus one per deeper level. With a 1\% false positive rate ($\phi = 0.01$): $T \times \phi + L \times \phi = 10 \times 0.01 + 3 \times 0.01 = 0.13$ expected false positive disk reads per lookup, compared to LevelDB's $L \times \phi = 0.04$. This $3.25\times$ increase is acceptable given the $10\times$ write amplification improvement.

% ============================================================
\section{Conclusion and Future Work}
\label{sec:conclusion}
% ============================================================

This paper presented KiwiDB, a complete LSM tree key-value store built from scratch in Python 3.12+. The system demonstrates that a modern, type-safe, async-native storage engine can be implemented from first principles while incorporating design innovations that address real limitations in existing C++ implementations --- specifically, native async I/O, parallel flush with formal ordering guarantees, and subprocess compaction for GIL bypass.

The engine implements the full LSM lifecycle: WAL-backed durable writes, concurrent skip list memtable, multi-level SSTable storage with bloom filters and sparse indexes, parallel flush pipeline with event-chain commit ordering, and hybrid tiering/leveling compaction with subprocess execution. A three-tier block cache, writer-preferred reader-writer locks, and backpressure-driven flow control complete the design. The accompanying REST API and web dashboard make the engine's internals fully observable in real time.

\subsection{Limitations}

KiwiDB deliberately trades feature breadth for implementation depth:

\begin{itemize}[leftmargin=*]
    \item \textbf{No range queries}: The engine supports only point lookups. The internal \texttt{KWayMergeIterator} could support range scans, but a public iterator API is not exposed.
    \item \textbf{No compression}: SSTables are stored uncompressed. Adding block-level compression (e.g., LZ4, zstd) to the SSTable writer would be a straightforward extension.
    \item \textbf{No transactions}: Only single-key atomicity is guaranteed. Multi-key \texttt{WriteBatch} semantics would require WAL group commit support.
    \item \textbf{Python performance ceiling}: Raw throughput is approximately $10{-}100\times$ lower than C++ implementations due to interpreter overhead.
\end{itemize}

\subsection{Future Work}

Several extensions are natural next steps:

\begin{itemize}[leftmargin=*]
    \item \textbf{Range query support}: Expose the internal merge iterator as a public API for ordered key-range scans across all levels.
    \item \textbf{Block compression}: Integrate LZ4 or zstd compression at the SSTable block level to reduce disk footprint and I/O bandwidth.
    \item \textbf{WAL group commit}: Batch multiple writes into a single fsync call to amortize the per-write fsync overhead under concurrent load.
    \item \textbf{Cross-level bloom optimization}: Allocate more bloom filter bits to deeper levels following the Monkey~\cite{dayan2017} approach to minimize total read amplification across levels.
    \item \textbf{GIL-free Python (PEP 703)}: When CPython's free-threaded mode matures, compaction could move from subprocess to thread pool, eliminating IPC overhead while retaining true parallelism.
\end{itemize}

% ============================================================
% References
% ============================================================

\begin{thebibliography}{00}

\bibitem{oneil1996}
P.~O'Neil, E.~Cheng, D.~Gawlick, and E.~O'Neil, ``The log-structured merge-tree (LSM-tree),'' \textit{Acta Informatica}, vol.~33, no.~4, pp.~351--385, 1996.

\bibitem{leveldb}
J.~Dean and S.~Ghemawat, ``LevelDB: A fast key-value storage library,'' Google, 2011. [Online]. Available: \url{https://github.com/google/leveldb}

\bibitem{rocksdb}
Facebook, ``RocksDB: A persistent key-value store for fast storage environments,'' 2013. [Online]. Available: \url{https://rocksdb.org}

\bibitem{cassandra}
A.~Lakshman and P.~Malik, ``Cassandra: A decentralized structured storage system,'' \textit{ACM SIGOPS Operating Systems Review}, vol.~44, no.~2, pp.~35--40, 2010.

\bibitem{wiredtiger}
MongoDB, ``WiredTiger storage engine,'' [Online]. Available: \url{https://source.wiredtiger.com}

\bibitem{bloom1970}
B.~H.~Bloom, ``Space/time trade-offs in hash coding with allowable errors,'' \textit{Communications of the ACM}, vol.~13, no.~7, pp.~422--426, 1970.

\bibitem{dayan2017}
N.~Dayan, M.~Athanassoulis, and S.~Idreos, ``Monkey: Optimal navigable key-value store,'' in \textit{Proc.\ ACM SIGMOD}, 2017, pp.~79--94.

\bibitem{dayan2018}
N.~Dayan and S.~Idreos, ``Dostoevsky: Better space-time trade-offs for LSM-tree based key-value stores via adaptive removal of superfluous merging,'' in \textit{Proc.\ ACM SIGMOD}, 2018, pp.~505--520.

\bibitem{herlihy2006}
M.~Herlihy, Y.~Lev, V.~Luchangco, and N.~Shavit, ``A provably correct scalable concurrent skip list,'' in \textit{Proc.\ OPODIS}, 2006.

\bibitem{luo2020}
C.~Luo and M.~J.~Carey, ``LSM-based storage techniques: A survey,'' \textit{The VLDB Journal}, vol.~29, pp.~393--418, 2020.

\end{thebibliography}

\end{document}
